%--beamer--------------------------------------------------------------------------------------
%	PACKAGES AND THEMES
%----------------------------------------------------------------------------------------
\documentclass[xcolor=dvipsnames]{beamer}
\usetheme{default}
\usepackage[backend=bibtex,
defernumbers=true,
style=numeric,
citestyle=ieee
]{biblatex}
\addbibresource{psref.bib}
\usepackage{hyperref}
\usepackage{graphicx} % Allows including images
\usepackage{booktabs} % Allows the use of \toprule, \midrule and \bottomrule in tables
\usepackage[utf8]{inputenc}
\usepackage[T1]{fontenc}
\usepackage{amsmath,amssymb}
\usepackage{multirow}
\usepackage{array}
\usepackage{pythontex}
\usepackage{hyperref}

% Uncomment here for a Turkish presentation.
%\AtBeginDocument{%
%	\renewcommand\tablename{Tablo}
%}
%\AtBeginDocument{%
%	\renewcommand\figurename{\u015eekil}
%}

%----------------------------------------------------------------------------------------
%	TITLE PAGE
%----------------------------------------------------------------------------------------

% The titlo
\title[Title]
{
 Quantum Algorithm for Roots of Multivariate Functions Over
 Finite Fields
}

\subtitle{Project Statement}
\author{ Alexander J. Heilman \& Andy Phillips }
 
% can specify the institute here
%\institute[]{\inst{1}Bilgisayar M�hendisli\u011fi B�l�m�\\Gebze Teknik �niversitesi\and\inst{2}Bili\u015fim Teknolojileri Enstit�s�\\Gebze Teknik �niversitesi}

\date{\today} % Date, can be changed to a custom date


%----------------------------------------------------------------------------------------
%	PRESENTATION SLIDES
%----------------------------------------------------------------------------------------

\begin{document}

\begin{frame}
    % Print the title page as the first slide
    \titlepage
\end{frame}

\begin{frame}{Overview}

$\bullet$ We can encode multivariate finite functions over
          finite fields in multi-qudit quantum 
          states \cite{appel2020}
          
\medskip\pause         
          
$\bullet$ Relatively new generalizations of Grover's search
          algorithm apply to multi-valued 
          functions \cite{hunt2020}\cite{fan2008}
          
\medskip\pause          
          
$\bullet$ Application of multi-valued search algorithms to 
          the known encoding should facilitate the evaluation 
          of roots
          
\medskip\pause

$\bullet$ Google's Cirq SDK allows simulation of qudit
          circuits

\end{frame}



\begin{frame}{Finite Function Encoding States}
Given an $n$-qu$d$it system, where $d$ is prime, we can encode an
$n$-variable finite function over the field $\mathbb{F}_d$ in the amplitude of the basis states.

$$
\vert f(x_n) \rangle \rightarrow \frac{1}{\sqrt{d^n}}\sum_{k=0}^{d^n-1} \omega_d^{f(k)}\vert k\rangle
$$
\end{frame}

\begin{frame}{Finite Function Encoding States Ex.}
Let's look at a very simple case of two qutrits (i.e. $d=3$). We can encode the polynomial function $xy^2+x$ over
the field $\mathbb{F}_d$ (i.e. the field with three elements, $0,1,2$) as follows:

$$
                   \frac{1}{\sqrt{9}}\sum_{k=0}^{9}\omega^{xy^2+x}\vert k \rangle        
$$

$$
      \hspace{-.5cm}    =\frac{1}{3}(\omega^{0\cdot0^2+0}\vert 00\rangle\pause + \omega^{0\cdot1^2+0}\vert 01\rangle\pause + \omega^{0\cdot2^2+0}\vert 02\rangle
          + \omega^{1\cdot0^2+1} \vert 10\rangle + \omega^{1\cdot1^2+1} \vert 11\rangle  
          $$
          \vspace{-.7cm}
          $$ \hspace{.83cm}+ \omega^{1\cdot2^2+1}\vert 12\rangle +\omega^{2\cdot0^2+2} \vert 20\rangle 
          + \omega^{2\cdot1^2+2} \vert 21\rangle + \omega^{2\cdot2^2+2} \vert 22\rangle  )   
        $$
        
        $$
          =\frac{1}{3}(\vert 00\rangle + \vert 01\rangle + \vert 02\rangle
          + \omega \vert 10\rangle + \omega^2 \vert 11\rangle +\hspace{.7cm}
          $$
          \vspace{-.63cm}
          $$ \hspace{1.7cm} \omega^2 \vert 12\rangle +\omega^2 \vert 20\rangle 
          + \omega \vert 21\rangle + \omega \vert 22\rangle  )   
        $$
        
\end{frame}        
        
\begin{frame}{$xy^2+x$ Example cont.}        
        
        $$\hspace{-.8cm}\vert xy^2+x\rangle = \begin{array}{c c}  x y
          & \begin{array}{c} \omega^{xy^2+x} \\ \end{array} \\
          \begin{array}{c c c c c c c c c}\vert 00\rangle \\ \vert 01 \rangle\\
          \vert 02 \rangle \\ \vert 10\rangle \\ \vert 11\rangle\\ \vert 12\rangle\\
          \vert 20 \rangle \\ \vert 21 \rangle \\ \vert 22\rangle\end{array} \hspace{-.5cm}&
          \left[
          \begin{array}{c}
          1 \\
          1 \\
          1 \\
          \omega \\
          \omega^2 \\
          \omega^2\\
          \omega^2\\
          \omega \\
          \omega \\
          \end{array}
          \right]
          \end{array}
        $$

\end{frame}


\begin{frame}{Multi-Valued Grover Search}
The typical Grover search algorithm is effectively used to find a 
set of basis states marked in a complete superposition of basis states by a relative amplitude of $-1$.

\vspace{1.2cm}\pause

A generalization of this is to find basis states marked with one of many relative amplitudes, which if 
equally spaced are the roots of unity $\omega_d=e^{2\pi i/d}$

\end{frame}



\begin{frame}{Multi-Valued Grover Search}
Maximilian Hunt and Samuel Hunt have recently published \textit{Grover’s Algorithm and Many-Valued Quantum Logic
}(December 2020, \cite{hunt2020}).\pause They generalize the Grover diffusion operator to qudits and multi-valued functions using the circuit below:

\includegraphics[scale=.45]{GGDO.pdf}

For gates see \url{https://alexheilman.com/qis/qudits.html}
\end{frame}


\begin{frame}{Goals/Expected Problems}

$\bullet$ Grover's search generally only works for sparse 
          databases\pause, those are collections
          where the solutions are a minority of the population
          \pause
          
          \medskip
          
$\bullet$ We should be able to at least count the number of
          roots using a generalized counting/amplitude
          estimation scheme
          

\end{frame}

\begin{frame}{Next Steps}

$\bullet$ Better understand multi-valued Grover search\pause

\medskip 

$\bullet$ Implement examples in Cirq simulations and get hands on\pause, specifically for the state from above and the 2-qutrit GGDO:

\includegraphics[scale=.45]{2qtGGDO.pdf}



\end{frame}





\begin{frame}[allowframebreaks]{References}
    % This might take more than one page
    \nocite{*}
    \printbibliography

\end{frame}


\end{document}
